\documentclass{beamer}
\usetheme{Madrid}
\usecolortheme{default}

\title[Register Allocation]{DA2026\_PRJ2: Compiler Register Allocation}
\subtitle{Graph-Coloring Heuristics}
\author[Group 4 - Class 5]{Beatriz Remondes \and Nuno Coimbra \and Simão Ribeiro}
\institute[FEUP]{Faculty of Engineering, University of Porto}
\date{Spring 2026}

\begin{document}

\begin{frame}
    \titlepage
\end{frame}

\begin{frame}{Outline}
    \tableofcontents
\end{frame}

\section{Introduction}
\begin{frame}{Project Overview}
    \begin{itemize}
        \item \textbf{Goal:} Implement a compiler back-end tool for global register allocation.
        \item \textbf{Approach:} Graph-coloring heuristics.
        \item \textbf{Pipeline:}
        \begin{enumerate}
            \item Parse live ranges from 3-address intermediate representation.
            \item Build \textbf{live webs} (merging overlapping ranges).
            \item Construct an \textbf{interference graph}.
            \item Assign webs to $N$ available registers.
        \end{enumerate}
        \item \textbf{Fallback Strategies:} Spilling, Splitting, and a Hybrid "Free" strategy when allocation fails.
    \end{itemize}
\end{frame}

\section{The Graph Class}
\begin{frame}{The Graph Class: Conceptual Decisions}
    \begin{itemize}
        \item \textbf{Base Structure:} Utilized the generic \texttt{Graph<T>} and \texttt{Vertex<T>} templates provided in practical classes.
        \item \textbf{Alterations \& Additions:}
        \begin{itemize}
            \item \textbf{Undirected Semantics:} Interference graphs are inherently undirected. Edge additions ensure symmetry ($A \rightarrow B$ and $B \rightarrow A$).
            \item \textbf{Node/Edge Removal:} Utilized \texttt{removeVertex} to dynamically simulate web spilling and graph simplification.
            \item \textbf{Incoming Edges Tracking:} Added \texttt{getIncoming()} and indegree management to support bidirectional edge analysis.
        \end{itemize}
    \end{itemize}
\end{frame}

\section{Implemented Heuristics}
\begin{frame}{Spilling Heuristic (T2.2)}
    \begin{itemize}
        \item \textbf{Objective:} Move uncolorable webs to memory to reduce graph chromatic number.
        \item \textbf{Spill Metric:} $\frac{\text{Degree(Node)}}{\text{Live\_Range\_length(Node)}}$
        \item \textbf{Rationale:} 
        \begin{itemize}
            \item High degree: removes maximum interference.
            \item Short live range: minimizes memory load/store overhead.
        \end{itemize}
        \item Iterates from $1$ to $K$ spilled webs, stopping when the graph becomes $N$-colorable.
    \end{itemize}
\end{frame}

\begin{frame}{Splitting Heuristic (T2.3)}
    \begin{itemize}
        \item \textbf{Objective:} Divide a web into two smaller webs to reduce localized interference.
        \item \textbf{Maximum Interference Bottleneck:}
        \begin{itemize}
            \item Analyzes internal execution points of the uncolorable web.
            \item Finds the specific program line with the highest simultaneous interference.
        \end{itemize}
        \item \textbf{Rationale:} Splitting at this exact point severs the most edges in the interference graph in a single operation.
    \end{itemize}
\end{frame}

\begin{frame}{"Free" Algorithm Strategy (T2.4)}
    \begin{itemize}
        \item \textbf{Dynamic Efficiency Threshold Approach.}
        \item Evaluates the web with the highest degree when an impasse is reached.
        \item \textbf{Split Efficiency Metric:} $\frac{\text{Live\_Range\_length(Node)}}{\text{Degree(Node)}}$
        \item \textbf{Adaptive Decision Making:}
        \begin{itemize}
            \item \textbf{Spill:} If efficiency $< 1.5 / N$ (dynamic cutoff) or web spans $\le 2$ lines.
            \item \textbf{Split:} If efficiency is above the threshold.
        \end{itemize}
        \item Prevents fragment pollution from over-splitting highly connected, short-lived webs.
    \end{itemize}
\end{frame}

\section{Challenging Aspects}
\begin{frame}{Challenging Aspects of Implementation}
    \begin{itemize}
        \item \textbf{Web Construction \& Merging:} 
        \begin{itemize}
            \item Handling continuous vs. discontinuous live ranges accurately.
            \item Merging variables with overlapping scopes efficiently $O(R^2 \cdot L)$.
        \end{itemize}
        \item \textbf{Splitting Logic:} 
        \begin{itemize}
            \item Ensuring semantic correctness when dividing a web (part 1 \& part 2).
            \item Efficiently rebuilding the interference graph after topological changes.
        \end{itemize}
        \item \textbf{Hybrid Strategy Tuning:} 
        \begin{itemize}
            \item Balancing the heuristics to correctly identify when to Spill vs. Split.
        \end{itemize}
    \end{itemize}
\end{frame}

\section{Demo}
\begin{frame}{Live Demonstration}
    \begin{center}
        \Large \textbf{Time for a Live Demo!}
    \end{center}
    \vspace{0.5cm}
    \begin{itemize}
        \item \textbf{Interactive Mode:} We will run our allocator using the step-by-step interactive menu.
        \item \textbf{Test Cases:} Resolving a sample dataset of live ranges and available registers.
    \end{itemize}
\end{frame}

\section{Conclusion}
\begin{frame}{Conclusion}
    \begin{itemize}
        \item Successfully implemented a robust register allocator using graph coloring.
        \item Extended the base \texttt{Graph} class for advanced heuristic support.
        \item The "Free" strategy optimally balances spills and splits.
    \end{itemize}
\end{frame}

\end{document}
